\Section{Действительные числа.
Теорема о существовании и единственности точной верхней (нижней)
грани числового множества, ограниченного сверху (снизу).
Счетность \(\Q\), несчетность \(\R\)}

\Subsection{Действительные числа}

\begin{definition}
\textit{Рациональным числом} называют дробь вида \(\dfrac{p}{q}\), где \(p\in\Z,~q\in\N\).
\end{definition}

\begin{definition}
\textit{Сечением множества рациональных чисел \(\Q\)} называют
разбиение этого множества на два непустых множества \(A_{*}\) и \(A^{*}\) таких, что:

\begin{itemize}
\item \(A_{*}\cup A^{*}=\Q\)
\item \(A_{*}\cap A^{*}=\emptyset\)
\item \(\forall\, x\in A_{*},~\forall\,y\in A^{*}\longmapsto x < y\)
\end{itemize}

При этом множества \(A_{*}\) и \(A^{*}\) называют
\textit{нижним} и \textit{верхним классами сечения} соответственно.
\end{definition}

\begin{notation}
Будем называть сечения, в которых существует наименьший элемент в \(A^{*}\),
но не существует наибольший элемент в \(A_{*}\), \textit{сечениями класса А}.

\vspace{2mm}

Если же, наоборот, отсутствует наименьший элемент в \(A^{*}\), но есть наибольший
элемент в \(A_{*}\), то будем относить такие сечения к \textit{классу В}.

\vspace{2mm}

В случае, когда нет ни наибольшего элемента в \(A_{*}\), ни наименьшего в \(A^{*}\),
будем относить такие сечения к \textit{классу С}.\\
\\
Поскольку классы \(A\) и \(B\) аналогичны, то для определенности будем
рассматривать только классы \(A\) и \(C\).
\end{notation}

\begin{definition}
\textit{Иррациональным числом} называется сечение класса \(C\). \\
\textit{Действительным числом} называются сечения любого из трех классов.
\end{definition}

\begin{definition}
\textit{Сечение множества действительных чисел \(\R\)} определяется
аналогично сечению \(\Q\).
\end{definition}

\begin{definition}
Множество \(X\subset\R\) называют \textit{ограниченным сверху},
если \(\exists\,\alpha\in\R\) такое, что для любого элемента \(x\in X\)
выполняется \(x\leq\alpha\). \\
При этом число \(\alpha\) называют \textit{верхней гранью} множества \(X\).
Ограниченность снизу и нижняя грань определяются аналогично.

\vspace{2mm}

Множество \(X\subset\R\) называют \textit{ограниченным},
если \(\exists\,\alpha\in\R\) такое, что для любого элемента \(x\in X\)
выполняется \(|x|\leq\alpha\). \\
Можно также сказать, что если множество ограничено сверху и снизу, то оно ограничено.
\end{definition}

\begin{definition}
Наименьшую из верхних граней множества \(X\) называют \textit{точной верхней гранью}.
\[
    [\alpha_{0} = \sup X] \overset{def}{=} [\forall \, x \in X \longmapsto \alpha_{0} \geq x]
    ~ \& ~ [\forall \, \alpha \in \R, \alpha < \alpha_{0} ~\exists \, x_{\alpha} \in X :
    x_{\alpha} > \alpha]
\]
Аналогично определяется \textit{точная нижняя грань}.
\[
    [\alpha_{0} = \inf X] \overset{def}{=} [\forall \, x \in X \longmapsto \alpha_{0} \leq x]
    ~ \& ~[\forall \, \alpha \in \R, \alpha > \alpha_{0} ~ \exists \, x_{\alpha} \in X :
    x_{\alpha} < \alpha]
\]
\end{definition}

\begin{lemma} \label{lemma_rational_in_real}
Для любых действительных чисел \(\alpha\) и \(\beta\) таких, что \(\alpha < \beta\), существует
\textbf{рациональное} число \(c\), для которого верно соотношение \(\alpha < c < \beta\).

\vspace{2mm}

\proofname~Пусть числам \(\alpha\) и \(\beta\) соответствуют сечения \textit{рациональных} чисел:
\[
\alpha = (A_{*}, \, A^{*}) \qquad \beta = (B_{*}, \, B^{*}).
\]
Поскольку \(\alpha < \beta\), то по определению операции сравнения \(A_{*} \subset B_{*}\)
и \(A_{*} \neq B_{*}\). Следовательно, найдется рациональное число \(p\) такое, что
\[
p \notin A_{*} \qquad \& \qquad p \in B_{*}.
\]
Поскольку множество \(B_{*}\) не имеет наибольшего элемента (напомним, что
мы рассматриваем только сечения класса \(A\)), то также найдется рациональное число \(p\)
такое, что
\[
q \notin A_{*} \qquad \& \qquad q \in B_{*} \qquad \& \qquad p < q.
\]

Тогда рассмотрим рациональное число \(r = \dfrac{p+q}{2}\). Заметим, что \(p < r < q\).
Осталось доказать, что \(\alpha < r < \beta\). \\

Из того, как мы определили числа \(p\) и \(q\) (и того, что мы используем сечения класса \(A\))
следует, что:
\[
\alpha < p \leq \beta,
\]
\[
\alpha < q \leq \beta.
\]
Следовательно, верно следующее:
\[
\alpha < p < r,
\]
\[
r < q \leq \beta.
\]
Окончательно:
\[
\alpha < r < \beta.
\]
Тогда мы доказали, что число \(r\) - искомое.~\qedsymbol
~\qedsymbol
\end{lemma}

\begin{lemma}
Не существует сечения действительных чисел такого,
что одновременно в нижнем классе сечения \(A_{*}\) содержится наибольший элемент,
а в верхнем классе сечения \(A^{*}\) содержится наименьший элемент.

\vspace{2mm}

\proofname~ Предположим противное: пусть \(\alpha_{0}\)~--- наименьшее число в \(A^{*}\), а
\(\beta_{0}\)~--- наибольшее число в \(A_{*}\). Заметим, что по определению сечений
\(\alpha_{0} > \beta_{0}\).

Тогда по Лемме 1.1 найдется такое действительное число \(c\), что
\(\alpha_{0} > c > \beta_{0}\). Но тогда число \(c\) не содержится ни в нижнем,
ни в верхнем классе сечений~--- следовательно, не выполняется первое свойство сечений:
\[
    A_{*}\cup A^{*} \neq \R
\]

Получили противоречие~--- следовательно, предположение неверно, и лемма верна.~\qedsymbol
\end{lemma}

\begin{theorem}[Теорема Дедекинда]
Для любого сечения действительных чисел либо существует наибольший элемент в \(A_{*}\),
либо существует наименьший элемент в \(A^{*}\).

\vspace{2mm}

\textit{Примечание: теорема не требует доказательства в рамках курса.}
\end{theorem}

\begin{theorem}
Любому действительному числу взаимно однозначно соответствует бесконечная десятичная дробь.

\vspace{2mm}

\textit{Примечание: теорема не требует доказательства в рамках курса.}
\end{theorem}

\begin{note}
Заметим, что \(0.9999\ldots=1.0000\ldots\). Тогда будем считать это одним и тем же числом,
и не рассматривать бесконечные дроби, заканчивающиеся на \(9999\ldots\).
\end{note}

\Subsection{Теорема о существовании и единственности точной верхней (нижней)
грани числового множества, ограниченного сверху (снизу)}

\begin{theorem}[Теорема о единственности точной верхней (нижней) грани]
Если непустое множество \(X\subset\R\) имеет точную верхнюю (нижнюю) грань, то эта грань~--- единственна.

\vspace{2mm}

\proofname~Докажем для точной верхней грани (для случая нижней~--- аналогично).\\
Предположим противное: пусть \(\alpha_{0}\) и \(\alpha'_{0}\)~---
точные верхние грани множества \(X\), причем \(\alpha_{0}\neq\alpha'_{0}\).

Без ограничения общности будем считать, что \(\alpha_{0}>\alpha'_{0}\).
Но тогда по определению точной верхней грани (для \(\alpha_{0}\)):
\[
\alpha'_{0} < \alpha_{0} ~ \Rightarrow ~ \exists\,x_{\alpha'_{0}}\in X:
\alpha'_{0} < x_{\alpha'_{0}}
\]

Следовательно, \(\alpha'_{0}\) не может являться точной верхней гранью множества \(X\).
Тогда наше предположение неверно, и точная верхняя грань единственна.~\qedsymbol
\end{theorem}

\begin{theorem}[Теорема о существовании точной верхней (нижней) грани] \label{theorem_set_SupInf}
Если множество \(X\subset\R\) ограничено сверху (снизу),
то оно имеет точную верхнюю (нижнюю) грань.

\vspace{2mm}

\proofname~Докажем для случая ограниченности сверху (для ограниченности снизу~--- аналогично). \\
Пусть \(X\subset\R\)~--- ограниченное сверху множество. Рассмотрим 2 случая: \\

\textbf{1.} Множество \(X\) имеет наибольший элемент. \\
Тогда этот элемент и будет точной верхней гранью: действительно,
любой элемент из \(X\) не больше него, и при этом любое число,
меньшее этого элемента, верхней гранью уже не является.\\

\textbf{2.} Множество \(X\) не имеет наибольшего элемента.\\
Построим такие 2 множества: множество \(A^{*}\) будет множеством верхних граней \(X\),
а множество \(A_{*}\) будет содержать все остальные действительные числа. Тогда:

\begin{itemize}
\item \(A_{*} \cup A^{*} = \R\)
\item \(A_{*} \cap A^{*} = \emptyset\)
\item \(\forall\,x\in A_{*},~\forall\,y\in A^{*} \longmapsto x < y\)
\end{itemize}

Последнее свойство верно потому, что если \(x\geq y\), то \(x\),
будучи больше верхней грани, будет сам являться верхней гранью,
то есть принадлежать \(A^{*}\), чего быть не может.\\

Тогда по определению наши 2 множества образуют сечение.
По теореме Дедекинда этому сечению соответствует некоторое число \(\alpha\).
Заметим, что если какой-то элемент \({x\in X}\) принадлежал бы \(A^{*}\),
то, являясь верхней гранью, он оказался бы наибольшим элементом \(X\),
чего по нашему предположению быть не может.
Следовательно, множество \(X\) целиком содержится в нижнем классе сечения \(A_{*}\).\\

Тогда, если бы \(\alpha\) было наибольшим числом в \(A_{*}\) (по т. Дедекинда),
то оно было бы верхней гранью \(X\), то есть одновременно принадлежало бы и \(A^{*}\),
чего быть не может.
Следовательно, (по т. Дедекинда) \(\alpha\)~--- наименьшее число в \(A^{*}\).
Тогда \(\alpha\)~--- наименьшая из верхних граней~--- то есть точная верхняя грань.~\qedsymbol
\end{theorem}

\Subsection{Счетность \(\Q\). Несчетность \(\R\)}

\begin{theorem}
Множество рациональных чисел \(\Q\)~--- счетно.

\vspace{2mm}

\proofname~Будем рассматривать только несократимые дроби,
поскольку сократимые лишь дублируют их. \\

Рассмотрим положительные дроби вида \(\dfrac{p}{q}\). Назовем \textit{высотой}
дроби величину \(H = p+q\) (например, высоте 5 соответствуют дроби
\(\dfrac{1}{4}, \dfrac{2}{3}, \dfrac{3}{2}, \dfrac{4}{1}\)).

Заметим, что каждой высоте соответствует \textit{конечное} число дробей.
Расположим дроби следующим образом: по возрастанию высоты,
для дробей одной высоты~--- по возрастанию величины:
\[
\underbrace{\frac{1}{1}}_{H=2},~ \underbrace{\frac{1}{2}, \frac{2}{1}}_{H=3},
~\underbrace{\frac{1}{3},\frac{3}{1}}_{H=4}\dots
\]

Перенумеруем все дроби, расположенные таким образом:
номеру \(n\) соответствует дробь \(a_{n}\).
Для того, чтобы охватить все множество рациональных чисел, зададим
следующий ряд чисел:
\[
0, a_{1}, -a_{1}, a_{2}, -a_{2}, \dots
\]

Перенумеруем числа в этом ряду и обозначим как \(b_{n}\).
Теперь мы взаимно однозначно сопоставили
каждое натуральное число \(n\) и рациональное число \(b_{n}\).
Следовательно, по определению \(\Q\)~--- счетно.~\qedsymbol
\end{theorem}

\begin{lemma}
Любое бесконечное подмножество счетного множества также счетно.

\vspace{2mm}

\proofname~Пусть \(A\)~--- счетное множество, а множество \(B\subset A\).
Тогда элементы множества \(A\) можно расположить в порядке возрастания номеров:
\[
A = \{a_{1}, a_{2}, a_{3}, a_{4}, \dots\}
\]

<<Пройдемся>> по множеству \(A\) в порядке возрастания номеров.
Выберем первый элемент, входящий в множество \(B\), и обозначим его \(a_{n_{1}}\).
Пойдем дальше и аналогично выберем \(a_{n_{2}},a_{n_{3}},\dots\)\\
Тогда мы взаимно однозначно сопоставили элементы множества \(B\) и натуральные числа:
\[
a_{n_{k}} ~ \longleftrightarrow ~ k
\]

Следовательно, множество \(B\)~--- счетно.
\end{lemma}

\begin{theorem}
Множество действительных чисел \(\R\) несчетно.

\vspace{2mm}

По Лемме 1.3 если какое-то подмножество \(\R\) несчетно, то и \(\R\) несчетно.
Тогда докажем, что отрезок \([0, 1]\) несчетен. \\

Предположим противное: пусть отрезок \([0, 1]\) счетен. Тогда каждому натуральному номеру \(n\)
соответствует действительное число, а, следовательно, бесконечная десятичная дробь:
\[
n \longleftrightarrow 0.a_{1n}a_{2n}a_{3n}a_{4n}\dots
\]

Зададим такое действительное число \(\gamma\):
\[
\gamma = 0.b_{1}b_{2}b_{3}b_{4}\dots,
\]
причем \(\forall\,n\in\N\hspace{3mm} b_{n}\neq a_{nn},~ b_{n} \neq 9\).

\definecolor{dark_green}{HTML}{006400}
\clearpage
\begin{table}[!t]
\centering
\begin{tabular}{| >{\centering\arraybackslash}m{2em}|
>{\raggedright\arraybackslash}m{5cm}|}
\hline
~ & \textcolor{dark_green}{0.25156737}\dots \\
\hline
1 & 0.\digitbox{1}3528395\dots \\
2 & 0.4\digitbox{9}828134\dots \\
3 & 0.18\digitbox{7}02956\dots \\
4 & 0.689\digitbox{3}2812\dots \\
5 & 0.8929\digitbox{1}405\dots \\
6 & 0.24500\digitbox{0}00\dots \\
7 & 0.812358\digitbox{2}4\dots \\
8 & 0.1111111\digitbox{1}\dots \\
\dots & \dots\dots\dots \\
\hline
\end{tabular}
\caption{Составление числа \(\gamma\) (сверху)}
\label{Таблица 1}
\end{table}

Заметим, что \(\gamma\) не равно ни одному из тех действительных чисел,
которым сопоставлен натуральный номер~--- следовательно, счетность отрезка \([0, 1]\)
нарушилась, то есть наше предположение неверно, и отрезок \([0, 1]\) несчетен.\\
Следовательно, несчетно и множество \(\R\).~\qedsymbol
\end{theorem}